
\documentclass[12pt, letterpaper,twocolumn]{article}
\usepackage[utf8]{inputenc}
\usepackage[english]{babel}
\usepackage{float}
\usepackage[total={6.5in, 9in}]{geometry}
\usepackage{graphicx}
\usepackage{longtable}
\usepackage{comment}
\usepackage{lscape}

\usepackage{amssymb,amsmath,amsthm,enumerate}
\usepackage[utf8]{inputenc}
\usepackage{array}
\usepackage[parfill]{parskip}
\usepackage{graphicx}
\usepackage{caption}
\usepackage{subcaption}
\usepackage{bm}
\usepackage{amsfonts,amscd}
\usepackage[]{units}
\usepackage{listings}
\usepackage{multicol}
\usepackage{multirow}
\usepackage{tcolorbox}
\usepackage{physics}
\usepackage{xcolor}
\usepackage{hyperref}

% Enable colored hyperlinks
\hypersetup{colorlinks=true}

% The following three lines are for crossmarks & checkmarks
\usepackage{pifont}% http://ctan.org/pkg/pifont
\newcommand{\cmark}{\ding{51}}%
\newcommand{\xmark}{\ding{55}}%


\usepackage{listings}
\usepackage{listingsutf8}
\usepackage{amsmath, amsthm, amssymb}
\lstset{ %
  language=python,
  basicstyle=\scriptsize,           % the size of the fonts that are used for the code
  numbers=none,
  numberstyle=\footnotesize,          % the size of the fonts that are used for the line-numbers
  numbersep=4pt,                  % how far the line-numbers are from the code
  backgroundcolor=\color{white},      % choose the background color. You must add \usepackage{color}
  breaklines=true,                % sets automatic line breaking
  breakatwhitespace=false,        % sets if automatic breaks should only happen at whitespace
  extendedchars=true,
  inputencoding=utf8, 
    tabsize=8,
    aboveskip=5pt,
    upquote=true,
    showstringspaces=false,
    frame=single,
    showtabs=false,
    showspaces=false,
    showstringspaces=false,
    keywordstyle=\color{blue}\ttfamily\bfseries,
    stringstyle=\color{red}\ttfamily,
    commentstyle=\color{violet}\ttfamily,
    morecomment=[l][\color{magenta}]{\#}, 
    literate={á}{{\'a}}1 {é}{{\'e}}1 {í}{{\'i}}1 {ó}{{\'o}}1 {ú}{{\'u}}1 {'}{{'}}1 {ñ}{{\~n}}1 {&}{{\&}}1 ,                                        
}

%\usepackage[superscript,biblabel]{cite}
\usepackage{algorithm2e}
\renewcommand{\thealgocf}{}

% Bibliography settings
%\usepackage[style=apa]{biblatex}
%\setbeamertemplate{bibliography item}{\insertbiblabel}
%\addbibresource{references/rq1.bib}
%\addbibresource{references/rq2.bib}
%\addbibresource{references/rq3.bib}
%\addbibresource{references/rq4.bib}
%\addbibresource{references/rq5.bib}
%\addbibresource{references/methodology.bib}
%\addbibresource{references/references.bib}
% Glossary entries
\usepackage[acronym]{glossaries}
\newacronym{ML}{ML}{machine learning}
\newacronym{HRI}{HRI}{human-robot interactions}
\newacronym{RNN}{RNN}{Recurrent Neural Network}
\newacronym{LSTM}{LSTM}{Long Short-Term Memory}


\theoremstyle{remark}
\newtheorem*{remark}{Remark}
\theoremstyle{definition}

\newcommand{\empy}[1]{{\color{darkorange}\emph{#1}}}
\newcommand{\empr}[1]{{\color{cardinalred}\emph{#1}}}
\newcommand{\examplebox}[2]{
\begin{tcolorbox}[colframe=darkcardinal,colback=boxgray,title=#1]
#2
\end{tcolorbox}}

%\usetheme{default} 
%\def \i  {\item}
\def \ai {\item[] \quad \arrowbullet}
\newcommand \si[1]{\item[] \quad \bulletcolor{#1}}
\def \wi {\item[] \quad $\ \phantom{\Rightarrow}\ $}
\def \bi {\begin{itemize}\item}
\def \ei {\end{itemize}}
\def \be {\begin{equation*}}
\def \ee {\end{equation*}}
\def \bie {$\displaystyle{}
\def \eie {{\ }$}}
\def \bsie {\small$\displaystyle{}
\def \esie {{\ }$}\normalsize\selectfont}
\def \bse {\small\begin{equation*}}
\def \ese {\end{equation*}\normalsize}
\def \bfe {\footnotesize\begin{equation*}}
\def \efe {\end{equation*}\normalsize}
\renewcommand \le[1] {\\ \medskip \lefteqn{\hspace{1cm}#1} \medskip}
\def \bex {\begin{example}}
\def \eex {\end{example}}
\def \bfig {\begin{figure}}
\def \efig {\end{figure}}
\def \btheo {\begin{theorem}}
\def \etheo {\end{theorem}}
\def \bc {\begin{columns}}
\def \ec {\end{columns}}
\def \btab {\begin{tabbing}}
\def \etab {\end{tabbing}\svneg\svneg}
\newcommand \col[1]{\column{#1\linewidth}}
\def\vneg  {\vspace{-5mm}}
\def\lvneg {\vspace{-10mm}}
\def\svneg {\vspace{-2mm}}
\def\tvneg {\vspace{-1mm}}
\def\vpos  {\vspace{5mm}}
\def\lvpos {\vspace{10mm}}
\def\svpos {\vspace{2mm}}
\def\tvpos {\vspace{1mm}}
\def\hneg  {\hspace{-5mm}}
\def\lhneg {\hspace{-10mm}}
\def\shneg {\hspace{-2mm}}
\def\thneg {\hspace{-1mm}}
\def\hpos  {\hspace{5mm}}
\def\lhpos {\hspace{10mm}}
\def\shpos {\hspace{2mm}}


%\logo{\includegraphics[height=0.4in]{images/univalle.png}}

% commands to relax beamer and subfig conflicts
% see here: https://tex.stackexchange.com/questions/426088/texlive-pretest-2018-beamer-and-subfig-collide
\makeatletter
\let\@@magyar@captionfix\relax
\makeatother

\title{Report 1: Characterizing and understanding security risks through Security-Aware Mutation Testing of security configuration in RESTful APIs}
%\subtitle{Subtitle Of Presentation}

\author{
    \large{Carlos Andres Delgado Saavedra}\\
    \footnotesize \href{mailto:carlos.andres.delgado@correounivalle.edu.co}{carlos.andres.delgado@correounivalle.edu.co} 
    Advisors:\\
    Jesus A. Aranda, Universidad del Valle \\
    Gills Perrouin, University of Namur,
    James Ortiz, Universite of Namur
    }

\usepackage{tikz}
\usetikzlibrary{shapes.geometric, arrows, positioning}

\begin{document}

\maketitle

\section{Introduction}

The focus of this report is the exploration and development of the \texttt{mutpy} library as part of a doctoral thesis. The structure of the \texttt{mutpy} library is as follows:

\begin{verbatim}
|-- mutpy
|   |-- codegen.py
|   |-- commandline.py
|   |-- controller.py
|   |-- coverage.py
|   |-- __init__.py
|   |-- operators
|   |   |-- arithmetic.py
|   |   |-- base.py
|   |   |-- decorator.py
|   |   |-- exception.py
|   |   |-- inheritance.py
|   |   |-- __init__.py
|   |   |-- logical.py
|   |   |-- loop.py
|   |   |__ misc.py
|   |-- templates
|   |   |-- base.html
|   |   |-- detail.html
|   |   |__ index.html
|   |-- termcolor.py
|   |-- test
|   |   |-- __init__.py
|   |   |-- test_commandline.py
|   |   |-- test_controller.py
|   |   |-- test_coverage.py
|   |   |-- test_operators.py
|   |   |-- test_runners.py
|   |   |-- test_utils.py
|   |   |-- test_views.py
|   |   |__ utils.py
|   |-- test_runners
|   |   |-- base.py
|   |   |-- __init__.py
|   |   |-- pytest_runner.py
|   |   |__ unittest_runner.py
|   |-- utils.py
|   |__ views.py
\end{verbatim}

The operators are specified in the \texttt{operators} folder, which receives the abstract syntax tree (AST).

\section{Installation of the library}

I created a fork of this library here: \url{https://github.com/cardel/mutpy}, I detected some problems about the compability, so I need to create a Docker container:

\begin{lstlisting}[language=make]
FROM python:3.9
COPY . usr/src/app
WORKDIR /usr/src/app
RUN apt-get update && apt-get install -y git
RUN pip install pip --upgrade
RUN pip install -r requirements.txt
RUN sh setup.sh
RUN python -m unittest discover tests
ENTRYPOINT sh tests.sh
\end{lstlisting}

And I needed to add some scripts:

Setup.sh

\begin{lstlisting}[language=bash]
git clone https://github.com/cardel/mutpy.git
cd mutpy && python setup.py install
\end{lstlisting} 

test.sh


\begin{lstlisting}[language=bash]
mut.py --target run.py --unit-test tests/test_app.py -m
\end{lstlisting} 

\section{Exploration of Flask Configurations}

The current configurations of Flask are loaded into the main variable as follows:

\begin{lstlisting}
print("Current Flask configurations:")
for key, value in app.config.items():
    print(f"{key}: {value}")
app.run(port=3000, debug=True)
\end{lstlisting}

Where \texttt{app} is the main variable. Basic Flask configurations include:

\begin{lstlisting}
DEBUG: False
TESTING: False
PROPAGATE_EXCEPTIONS: None
SECRET_KEY: None
PERMANENT_SESSION_LIFETIME: 31 days, 0:00:00
USE_X_SENDFILE: False
SERVER_NAME: None
APPLICATION_ROOT: /
SESSION_COOKIE_NAME: session
SESSION_COOKIE_DOMAIN: None
SESSION_COOKIE_PATH: None
SESSION_COOKIE_HTTPONLY: True
SESSION_COOKIE_SECURE: False
SESSION_COOKIE_SAMESITE: None
SESSION_REFRESH_EACH_REQUEST: True
MAX_CONTENT_LENGTH: None
SEND_FILE_MAX_AGE_DEFAULT: None
TRAP_BAD_REQUEST_ERRORS: None
TRAP_HTTP_EXCEPTIONS: False
EXPLAIN_TEMPLATE_LOADING: False
PREFERRED_URL_SCHEME: http
TEMPLATES_AUTO_RELOAD: None
MAX_COOKIE_SIZE: 4093
\end{lstlisting}

Default security settings might include:

\begin{itemize}
    \item \texttt{SESSION\_COOKIE\_SECURE: True}
    \item \texttt{SESSION\_COOKIE\_SAMESITE: 'Strict'}
    \item \texttt{CSRF\_ENABLED: True}
    \item \texttt{X\_FRAME\_OPTIONS: 'SAMEORIGIN'}
\end{itemize}

\section{Development Environment and Docker Setup}

Due to compatibility issues, a Docker container was necessary. The setup is as follows:

\begin{lstlisting}
FROM python:3.9
COPY . /usr/src/app
WORKDIR /usr/src/app
RUN apt-get update && apt-get install -y git
RUN pip install pip --upgrade
RUN pip install -r requirements.txt
RUN sh script.sh
RUN python -m unittest discover tests
ENTRYPOINT sh tests.sh
\end{lstlisting}

The \texttt{script.sh} contains:

\begin{lstlisting}
#!/bin/bash
git clone https://github.com/cardel/mutpy.git
cd mutpy && python setup.py install
\end{lstlisting}

\texttt{tests.sh} runs:

\begin{lstlisting}
mut.py --target run.py --unit-test tests/test_app.py -m
\end{lstlisting}

\section{Mutation Operators}

Mutation operators interact with the AST of the file to be mutated, generating a mutant. For example:

\begin{lstlisting}
def mutate_FloorDiv_to_Div(self, node):
    if self.should_mutate(node):
        return ast.Div()
    raise MutationResign()

def mutate_FloorDiv_to_Mult(self, node):
    if self.should_mutate(node):
        return ast.Mult()
    raise MutationResign()
\end{lstlisting}


An example for disabling CORS:

\begin{lstlisting}
app.config["ALLOW HEADERS"] = "*"
\end{lstlisting}

The exploration and development with \texttt{mutpy} and modification to Flask applications are integral parts of this doctoral thesis progress.

\end{document}
